\chapter{Cơ sở lý thuyết}\label{ch:background}

\section{Software Supply Chain và các rủi ro bảo mật}

\subsection{Định nghĩa chuỗi cung ứng phần mềm}
% TODO: Định nghĩa chuỗi cung ứng phần mềm
% Source Code → Dependencies → Build Process → Artifacts → Registry → Deployment

\subsection{Các vector tấn công phổ biến}
% TODO: Phân tích các vector tấn công ở từng giai đoạn
% - Source: unauthorized commit, compromised repo
% - Build: compromised build process, cache poisoning
% - Dependency: dependency confusion, typosquatting
% - Distribution: registry poisoning, MITM

\subsection{Các cuộc tấn công tiêu biểu}
% TODO: Case studies
% - SolarWinds (2020): build process compromise
% - Codecov (2021): CI script modification
% - Log4Shell (2021): dependency vulnerability
% - xz-utils (2024): long-term social engineering

\section{Mô hình bảo mật Zero Trust}

\subsection{Khái niệm Zero Trust}
% TODO: Giới thiệu Zero Trust Architecture
% - Nguyên tắc "Never trust, always verify"
% - NIST SP 800-207

\subsection{Áp dụng Zero Trust vào SDLC}
% TODO: Cách áp dụng Zero Trust vào quy trình phát triển phần mềm
% - Không tin tưởng bất kỳ code commit nào chưa verified
% - Không tin tưởng build runner chưa xác thực
% - Không tin tưởng container image chưa có chữ ký hợp lệ
% - Verify at every stage

\section{Khung tiêu chuẩn SLSA}

\subsection{Tổng quan về SLSA}
% TODO: Giới thiệu SLSA
% - Nguồn gốc: Google Binary Authorization for Borg
% - OpenSSF project
% - Maturity model, không phải tool

\subsection{Build Track -- Các cấp độ bảo mật}
% TODO: Bảng phân tích Build L0 → L3
% L0: No guarantees
% L1: Automated build + provenance documentation
% L2: Hosted build platform + signed provenance  
% L3: Isolated/ephemeral build + unforgeable provenance

\subsection{Source Track}
% TODO: Source L1 → L4 (theo draft v1.2)
% Tập trung giải thích L1-L3 (phần liên quan đến đồ án)

\subsection{Threat Model theo SLSA}
% TODO: Phân tích mô hình đe dọa
% - Source threats
% - Build threats  
% - Packaging & Distribution threats
% Ánh xạ với các cấp độ SLSA

\subsection{Mô hình xây dựng và Provenance Schema}
% TODO: 
% - Build Model: Control Plane vs Build Environment
% - In-toto attestation format
% - Provenance schema: buildDefinition, runDetails
% - externalParameters, resolvedDependencies

\section{Kubernetes và Container Security}

\subsection{Kiến trúc Kubernetes}
% TODO: Tổng quan K8s
% - Control Plane: API Server, etcd, Scheduler, Controller Manager
% - Worker Nodes: kubelet, container runtime
% - Pods, Deployments, Services

\subsection{Admission Controllers trong Kubernetes}
% TODO: Giải thích cơ chế Admission Controller
% - Validating Webhooks
% - Mutating Webhooks
% - Luồng xử lý: API Request → Authentication → Authorization → Admission → etcd

\subsection{OCI Image Format và Image Security}
% TODO:
% - OCI image spec
% - Image digest vs tag (tại sao tag không an toàn)
% - Image signing concepts

\section{Hệ sinh thái công cụ liên quan}

\subsection{Sigstore -- Hệ sinh thái ký số}
% TODO: Giới thiệu Sigstore
% - Cosign: ký và xác thực container image
% - Fulcio: certificate authority (keyless signing via OIDC)
% - Rekor: transparency log
% - Keyless signing workflow

\subsection{SBOM -- Software Bill of Materials}
% TODO:
% - Định nghĩa SBOM
% - Các chuẩn: SPDX, CycloneDX
% - Công cụ: Syft, Trivy
% - Vai trò trong supply chain security

\subsection{Policy Engines}
% TODO: So sánh các policy engine
% - Kyverno: YAML-native, dễ sử dụng, tích hợp Cosign
% - OPA Gatekeeper: Rego language, mạnh nhưng phức tạp
% - Sigstore Policy Controller
% Giải thích tại sao chọn phương án linh hoạt
